\section{The original one-point choice and its zero-dimensional obstruction bound}
\label{sec:count-choice-v036}

\status{CCH1--CCH17 construct the complete original one-point
\texttt{Choice} output after the \texttt{NiceI} normalization, with its
full block $a=n$.  In the selected Count-II application $n=4$ and
$\ell=3$ is allowed.  The construction chooses a literal
$g=0$ \texttt{Kap}, proves its full twisted \texttt{Detects} condition,
and preserves it through one finite local coefficient extension.
The same chosen \texttt{Kap} then satisfies the literal \texttt{B2}
with $C=0$ at every truncation and every level.  The original
ordinary/control, four further Hida endpoints, same-ideal
\texttt{HasAut}, first-order bound and \texttt{Raw} remain inputs.}

Fix a prime $\ell\ne2$, an original $F:\texttt{GFrame}(\ell)$, and
$x:\texttt{NP}(F,n)$ with $x.a=n$.  Write
\[
 \begin{gathered}
 G=F.G_0=\ker\eta,\qquad H=F.G'=\ker(\epsilon|_G),\\
 k=x.K,\quad A=k[[t]],\quad E=k((t)),\quad
 r(g)=\operatorname{mat}(x.\rho(g)).
 \end{gathered}
 \tag{CCH.1}
\]
The prime on $G'$ in the source means the full cyclotomic kernel.
It does not mean the kernel of its residual character.  Let
$\sigma(g)=cgc^{-1}$ and let $\chi(g)$ be the image of
$\epsilon(g)$ in $A^\times$ and then in $E^\times$.
We also use $\chi(c)$ for the image of $\epsilon(c)$ in these rings,
and put $s=-\chi(c)$.  The hypotheses needed for choice are
\[
 \Phi_\ell(F.\Gamma),\qquad x.N3,\qquad
 \exists g_0\in G:\overline{\epsilon(g_0)}\ne1.
 \tag{CCH.2}
\]
Here $\Phi_\ell$ is the original \texttt{PhiCond}, with every open
subgroup and every open-kernel $\mathbb F_\ell$-valued character.
CCH1--CCH6 use only the indicated parts of these hypotheses;
CCH7--CCH8 are where $\Phi_\ell$ is needed.
All local groups, ordered characters, flags and the original
auxiliary subgroup $D_0$ are retained throughout.

% BEGIN CCH1
\paragraph{CCH1. The actual scalar, group and full-block identities.}
\label{claim:CCH1}
The rank satisfies $n\ge1$, by $x.\texttt{a\_pos}$ and $x.a=n$.
The field $k$ has characteristic $\ell$.  Indeed, if its prime
characteristic were $p\ne\ell$, then $p$ would be a unit in
$\mathbb Z_\ell$ whose image in $k$ is zero, which a unital ring
map cannot do.  Every map from $\mathbb Z_\ell$ to a ring of
characteristic $\ell$ factors through $\mathbb F_\ell$: write any
$z\in\mathbb Z_\ell$ as $a+\ell b$, with $a$ an integer representative
of its residue.  This also proves uniqueness of that map.  Thus
$\chi$ is the original residual cyclotomic character embedded as
constant power series; its entries are locally constant, and
\[
 \chi|_H=1,\quad \chi(g_0)\ne1,\quad
 \chi(\sigma(g))=\chi(g),\quad s^2=1.
 \tag{CCH.3}
\]
The last identity follows from $c^2=1$.  The group $H$ is normal
and contains every commutator of $G$, since it is the kernel of a
homomorphism to an abelian group.  The group $G$ is compact and
totally disconnected: it is the closed open subgroup $\ker\eta$
of the original profinite group.

In the literal \texttt{blkMat} definition, both indices of an
$n\times n$ matrix are less than $a=n$.  Consequently
\[
 \operatorname{Ev}(A,A,n,n)=M_n(A),\qquad
 \operatorname{Od}(A,A,n,n)=0.                            \tag{CCH.4}
\]
The \texttt{NP.span} field says that each matrix unit, after
multiplication by a power of $t$, is in the $A$-span of $r(H)$.
Inverting $t$ and then spanning the matrix units gives
\[
                         \operatorname{span}_E r(H)=M_n(E).
 \tag{CCH.5}
\]
The \texttt{NP.form} field, whose form matrix is exactly one, is
\[
 r(g)^T r(\sigma(g))=\chi(g)^{1-n}I,
 \qquad
 r(\sigma(g))=\chi(g)^{1-n}r(g)^{-T}.                   \tag{CCH.6}
\]
These are the original source identities.  No semisimplicity,
characteristic-zero hypothesis or division by $n!$ has been added.
\hfill$\square$
% END CCH1

% BEGIN CCH2
\paragraph{CCH2. Construct the original \texttt{Kap} with $g=0$.}
\label{claim:CCH2}
Define $k_0:\texttt{Kap}(x)$ by $k_0.g=0$, and take the unique
empty functions for $\kappa$ and its local corrections $Y$.
The \texttt{mem}, \texttt{Y\_mem} and \texttt{loc} fields are
empty universal statements.  For \texttt{indep}, a function
$\mathrm{Fin}(0)\to A$ is already the zero function.

For \texttt{gen}, let $z$ belong to the original
$\operatorname{SelZ}$ with values in $\operatorname{Od}$.
Its first field and (CCH.4) give $z(g)=0$ for every $g$.
Choose exponent zero and the empty coefficient function.  The
resulting zero cocycle is in \texttt{Bd}: choose the zero matrix,
which lies in $\operatorname{Od}$, has transpose its negative,
and has zero coboundary.  This verifies every \texttt{Kap} field,
with all original local groups and flags unchanged.

This construction is available for every full-block $x$; it
does not invoke the general source theorem \texttt{kapEx}.
The integer $g$ in \texttt{Kap} is the number of odd Selmer
generators, not the rank $n$ of the representation.
\hfill$\square$
% END CCH2

% BEGIN CCH3
\paragraph{CCH3. Identify the full original twisted detector.}
\label{claim:CCH3}
For $k_0.g=0$, the original coefficient ring
$\operatorname{DD}(A,0)$ has two nonempty coordinates,
the constant coordinate and the $\eta$ coordinate.  Write
\[
 C(a)=(a,0),\qquad \nu(a)=(0,a),\qquad
 (a,b)(a',b')=(aa',ab'+ba').                            \tag{CCH.7}
\]
This is its actual multiplication formula; the two sums indexed
by $\mathrm{Fin}(0)$ vanish.  In particular $\nu(A)^2=0$.
The \texttt{pairedDualFirstOrderDualMatrixSpace} at $a=n$
is exactly $\nu(M_n(A))$.  The original $\rho_1$ and
$\epsilon_D$ are $C(r(g))$ and $C(\chi(g))$.

It follows by taking the $\eta$ coordinate that
$Z\in\texttt{TwZ}(k_0,A)$ is equivalent to $Z(g)=\nu(z(g))$
for a function $z:G\to M_n(A)$ satisfying all three conditions
\[
 \begin{split}
 &\text{each entry of $z$ is locally constant modulo $t^N$ for every $N$},\\
 &z(gh)=z(g)+\chi(g)r(g)z(h)r(g)^{-1},\\
 &z(\sigma(g))=s\,z(g)^T.
 \end{split}
 \tag{CCH.8}
\]
For the topology statement, multiplication by $C(t^N)$ in
\texttt{DD} multiplies both coordinates by $t^N$; the quotient
test on $\nu(z_{ij})$ is therefore precisely the quotient test
on $z_{ij}$.  Thus the equivalence includes the full
\texttt{AdicCont} field, not just the algebraic cocycle law.

Likewise $Z\in\texttt{TwB}(k_0,A)$ means that for some
$Y\in M_n(A)$ with $sY^T=Y$,
\[
 Z(g)=\nu(d_\chi Y(g)),\qquad
 d_\chi Y(g)=\chi(g)r(g)Yr(g)^{-1}-Y.                 \tag{CCH.9}
\]
Finally the source's trace test is exactly
\[
            \operatorname{ell}_{k_0}(h,Z)
                    =\operatorname{tr}(z(h)r(h)).      \tag{CCH.10}
\]
Indeed $\nu(z)C(r)=\nu(zr)$, and trace commutes with taking
the $\eta$ coordinate.  Hence the required \texttt{Detects}
statement on a set $S\subseteq G$ is: vanishing of (CCH.10)
on $S$ forces $t^b z=d_\chi Y$ for an integral $s$-symmetric
$Y$ and some $b\ge0$.  Both the cyclotomic twist and the
transpose condition remain present.
\hfill$\square$
% END CCH3

% BEGIN CCH4
\paragraph{CCH4. Trace detection over the fraction field.}
\label{claim:CCH4}
Let $z:G\to M_n(E)$ satisfy the cocycle law in (CCH.8), and
assume $\operatorname{tr}(z(h)r(h))=0$ for all $h\in H$.
Then $z=d_\chi Y$ for some $Y\in M_n(E)$.

\emph{Proof.}
On $H$ the twist is one.  Put $F(h)=z(h)r(h)$, so
\[
 F(hh')=F(h)r(h')+r(h)F(h'),\qquad \operatorname{tr}F(h)=0.
 \tag{CCH.11}
\]
Suppose $\sum_h a_h r(h)=0$, with finitely many nonzero $a_h$.
For every $h'\in H$, (CCH.11) gives
\[
 \operatorname{tr}\left(\left(\sum_h a_hF(h)\right)r(h')\right)
 =\sum_h a_h\operatorname{tr}F(hh')
       -\operatorname{tr}\left(\left(\sum_h a_hr(h)\right)F(h')\right)=0.
\]
By (CCH.5) and the nondegenerate matrix trace pairing,
$\sum_h a_hF(h)=0$.  The pairing is nondegenerate in every
characteristic: testing a matrix against $E_{ji}$ reads its
$(i,j)$ entry.  Therefore $r(h)\mapsto F(h)$ extends to an
$E$-linear map $D:M_n(E)\to M_n(E)$.  Equation (CCH.11)
and bilinearity show that $D$ is a derivation.

Here is its inner matrix explicitly.  Using indices
$1,\ldots,n$, set $Q=\sum_iD(E_{i1})E_{1i}$.
Differentiate $E_{ab}E_{i1}=\delta_{bi}E_{a1}$ and multiply
on the right by $E_{1i}$, then sum over $i$.  Since
$\sum_iE_{i1}E_{1i}=I$, this gives
\[
             QE_{ab}-E_{ab}Q=D(E_{ab}).                \tag{CCH.12}
\]
Thus $D(M)=QM-MQ$ for all matrices.  On $H$ we have
$z(h)=Q-r(h)Qr(h)^{-1}=d_\chi(-Q)(h)$.
Subtract this coboundary and denote the remaining cocycle by $z_0$;
it vanishes on $H$.

For $h\in H$ and $g\in G$, compare $hg$ with
$g(g^{-1}hg)$.  Normality of $H$ and the cocycle law show
$r(h)z_0(g)r(h)^{-1}=z_0(g)$.  By (CCH.5), $z_0(g)$
commutes with every matrix, and hence is scalar.  This last
assertion follows directly by commuting with all $E_{ii}$ and
$E_{ij}$.  Consequently
\[
 z_0(gh)=z_0(g)+\chi(g)z_0(h).
\]
Since $[G,G]\subseteq H$, the same formula gives
$z_0(gh)=z_0(hg)$.  Choose the $g_0$ of (CCH.3); comparison
of these two products yields
\[
 (\chi(g_0)-1)z_0(g)=(\chi(g)-1)z_0(g_0).
\]
The scalar matrix $Q_0=(\chi(g_0)-1)^{-1}z_0(g_0)$ therefore
satisfies $d_\chi Q_0=z_0$.  Taking $Y=-Q+Q_0$ proves the claim.
\hfill$\square$
% END CCH4

% BEGIN CCH5
\paragraph{CCH5. Polarized integral detection.}
\label{claim:CCH5}
The original $\texttt{Detects}(k_0,H)$ holds.
Take an integral $z$ as in CCH3 and apply CCH4 to obtain
$z=d_\chi Y$ over $E$.  Equation (CCH.6), with its scalar
canceling from conjugation, gives the exact identity
\[
             (d_\chi Y(\sigma(g)))^T=d_\chi(Y^T)(g).
 \tag{CCH.13}
\]
The polarization of $z$ implies $d_\chi(Y^T)=s z$.
Because $2$ is invertible and $s^2=1$, the matrix
\[
             Y_s=\tfrac12(Y+sY^T)
\]
satisfies $sY_s^T=Y_s$ and $d_\chi Y_s=z$.
Choose a single $b\ge0$ clearing all Laurent-series
denominators of its finitely many entries.  Then
$Y_A=t^bY_s\in M_n(A)$, still $sY_A^T=Y_A$, and
$t^b z=d_\chi Y_A$.  By CCH3 this is exactly the existential
power and original \texttt{TwB} required by \texttt{Detects}.
Only $2$ and the nonzero scalar $\chi(g_0)-1$ were inverted
over $E$; in particular $\ell=3,n=4$ causes no exception.
\hfill$\square$
% END CCH5

% BEGIN CCH6
\paragraph{CCH6. Detection on the actual separable subset.}
\label{claim:CCH6}
Set
\[
 S=\{h\in H:\det(XI-r(h))\text{ is separable over }E\}.
 \tag{CCH.14}
\]
The literal $x.N3$ says precisely that every open set meeting
$H$ meets $S$.  Hence $H\subseteq\overline S$.
For an integral cocycle $z$ in (CCH.8), the function
$h\mapsto\operatorname{tr}(z(h)r(h))$ is locally constant
modulo each $t^N$, by finite addition and multiplication of
the original entrywise quotient-continuous functions.
Its zero set is the intersection of the closed zero sets
modulo $t^N$, because $\bigcap_Nt^NA=0$.
Vanishing on $S$ therefore implies vanishing on $H$.
CCH5 proves
\[
                         \texttt{Detects}(k_0,S).
 \tag{CCH.15}
\]
This one-point argument uses the exact $N3$ already supplied
by CIP10.  It does not require the additional simultaneous
separability argument for a family of several points.
\hfill$\square$
% END CCH6

% BEGIN CCH7
\paragraph{CCH7. Finite generation of the actual continuous twisted cocycles.}
\label{claim:CCH7}
Let $\mathcal Z$ be the $A$-module of all entrywise
\texttt{AdicCont} functions $z:G\to M_n(A)$ with
\[
 z(gh)=z(g)+\chi(g)r(g)z(h)r(g)^{-1}.
 \tag{CCH.16}
\]
We do not impose polarization in this larger module.
Under $\Phi_\ell(F.\Gamma)$, $\mathcal Z$ is finitely generated
over $A$.

\emph{Proof.}
Put $M=M_n(A)$ and $\overline M=M/tM$.
The action $g\cdot Z=\chi(g)r(g)Zr(g)^{-1}$ is $A$-linear.
Modulo $t$, it is a continuous action on the finite additive
$\mathbb F_\ell$-vector space $\overline M$.  Its kernel
$U\subseteq G$ is open and normal.  It is also open in
$F.\Gamma$, because $G$ is open.  The finite quotient $G/U$
has a finite set of representatives, by compactness.

There are finitely many continuous $1$-cocycles
$G\to\overline M$.  To see this directly, restrict one to
$U$, where it is a continuous additive homomorphism.
Choosing an $\mathbb F_\ell$-basis of $\overline M$ identifies
such homomorphisms with a finite product of the original
$\operatorname{Hom}_\ell(U)$ sets, each finite by
$\Phi_\ell$.  A cocycle is then determined by this restriction
and its values on the finitely many coset representatives:
$z(gu)=z(g)+g\cdot z(u)$.  There are only finitely many
choices for these data.

Reduction modulo $t$ maps $\mathcal Z$ into this finite set.
Its kernel is exactly $t\mathcal Z$.  The nontrivial direction
is division by $t$: if every entry of $z(g)$ is divisible by
$t$, define $w(g)=z(g)/t$.  The torsion-free matrix module
allows cancellation of $t$ in the cocycle equation.  The
reduction of $w$ modulo $t^N$ is determined by that of $z$
modulo $t^{N+1}$, so $w$ has the same required adic continuity.
Thus $\mathcal Z/t\mathcal Z$ is a finite-dimensional
$k$-vector space.

Choose lifts $z_1,\ldots,z_d$ of a basis of that quotient.
For $z\in\mathcal Z$, recursively subtract a $k$-linear
combination of these lifts and divide the remainder by $t$.
After $N$ steps,
\[
 z=\sum_{i=1}^d a_{i,N}z_i+t^Nw_N,
 \qquad a_{i,N}\in k[t],\quad \deg a_{i,N}<N,
 \tag{CCH.17}
\]
where successive coefficient polynomials agree modulo $t^N$.
They define $a_i\in k[[t]]$.
For every $N$, $z-\sum_i a_i z_i\in t^N\mathcal Z$:
subtract the tails $(a_i-a_{i,N})/t^N$ times the generators
from $w_N$ in (CCH.17).  Evaluating at any $g$, each entry
lies in every $t^NA$, and is zero.  Therefore
$z=\sum_i a_i z_i$.  If $d=0$ the same argument gives
$\mathcal Z=0$.

This proves finite generation by an explicit power-series
lifting argument.  No replacement of $\Phi_\ell$ by a
topological finite-generation hypothesis has occurred.
\hfill$\square$
% END CCH7

% BEGIN CCH8
\paragraph{CCH8. A finite detector consisting of the required separable elements.}
\label{claim:CCH8}
For $h\in S$ define the $A$-linear functional
$L_h(z)=\operatorname{tr}(z(h)r(h))$ on $\mathcal Z$.
Choose the generators $z_1,\ldots,z_d$ from CCH7, and form
the rows
\[
                       v_h=(L_h(z_1),\ldots,L_h(z_d))\in E^d.
\]
Choose finitely many of these rows forming a basis of the
$E$-span of all of them.  Such a choice needs at most $d$
rows: add a row outside the existing span until its dimension
can no longer increase.  Let $S_0\subseteq S$ be the associated
finite set of elements.

For $z=\sum_i a_i z_i$, vanishing of $L_h(z)$ for $h\in S_0$
implies vanishing for every $h\in S$, by expressing $v_h$ as
an $E$-linear combination of the selected rows and multiplying
by the coefficient column $(a_i)$.  Its value lies in $A$,
whose embedding in $E$ is injective.  CCH6 then yields
\[
                         \texttt{Detects}(k_0,S_0).
 \tag{CCH.18}
\]
Enumerate $S_0$ by $\mathrm{Fin}(q)$ and call the resulting
map $\mathrm{sg}$.  Every $\mathrm{sg}(a)$ lies in the
ORIGINAL $G'$ and has separable characteristic polynomial
over $E$.  The detector is the complete predicate of CCH3.
The construction allows $q=0$ when the finite tests are empty;
the source also allows that value.
\hfill$\square$
% END CCH8

% BEGIN CCH9
\paragraph{CCH9. The needed finite complete DVR extension, with proof.}
\label{claim:CCH9}
Let $L/E$ be any finite separable extension, and let $B$ be
the integral closure of $A=k[[t]]$ in $L$.  Then $B$ is a
finite free $A$-module of positive rank, a complete discrete
valuation ring with fraction field $L$, and has finite residue
field.  The inclusion $A\hookrightarrow B$ is injective and local.

\emph{Proof.}
The first nonzero coefficient defines a discrete valuation on
$A$; every nonzero ideal is generated by a power of $t$.
Thus $A$ is a principal ideal domain and is integrally closed:
an element of $E$ with negative valuation cannot satisfy a
monic equation over $A$, since the leading power has strictly
smaller valuation than all other terms.

The trace pairing of $L/E$ is nondegenerate.  For clarity,
$E=k((t))$ is infinite, and a finite separable extension has a
primitive element: for two generators, exclude the finitely
many scalars that make two conjugates of their linear
combination equal, and then iterate.  On the power basis of
a primitive element the trace Gram matrix is $V^TV$, with
$V$ the Vandermonde matrix of its distinct conjugates.
Its determinant is nonzero, also in characteristic $\ell$.

Choose an $E$-basis $b_1,\ldots,b_d$ of integral elements
by multiplying any basis by suitable powers of $t$.
This scaling is possible by clearing the coefficients in
their monic equations over $E$.  The trace of an integral
element belongs to $A$: its conjugates are integral, their
sum belongs to $E$, and $A$ is integrally closed.
If $b_i^*$ is the trace-dual basis, every $y\in B$ satisfies
\[
                  y=\sum_i\operatorname{Tr}_{L/E}(yb_i)b_i^*
                        \in\sum_i A b_i^*.             \tag{CCH.19}
\]
It follows that $B$ is finite over $A$.  Here a submodule of
a finite free module over a principal ideal domain is finite
free: project to one coordinate, choose a generator of the
image ideal and a lift, split off that generator, and induct
on the number of coordinates.  Applied to (CCH.19), this
also proves freeness.  Its rank is positive, and $B$ spans
$L$, since it contains the chosen basis.

Hence $B$ is complete and separated for the $t$-adic topology,
and $B/tB$ is a finite ring.  Also $t$ is not a unit in $B$:
otherwise $t^{-1}\in E$ would be integral over $A$.
Every maximal ideal of $B$ contains $t$.  Indeed $B$ is integral
over $A$, and if a field is integral over a subdomain then the
subdomain is a field: the monic equation for $a^{-1}$, multiplied
by a suitable power of $a$, puts $a^{-1}$ in the subdomain.
Apply this to $A/(\mathfrak m\cap A)\subseteq B/\mathfrak m$.

We next prove that $B$ is local.  A finite-dimensional
commutative $k$-algebra has nilpotent Jacobson radical.  Its
radical powers stabilize, and Nakayama's elementary determinant
argument makes the stable power zero: if generators of a
module satisfy $v=Cv$ with all entries of $C$ in the radical,
then $\det(I-C)$ is a unit annihilating them.
The Chinese remainder theorem for powers of its distinct
maximal ideals consequently gives a product of local rings.
In particular, more than one maximal ideal would give a
nontrivial idempotent in $B/tB$.
Any such idempotent lifts to $B$.  Starting with a lift $e$,
the Newton step
\[
            e\longmapsto e-(e^2-e)(2e-1)^{-1}
\]
squares the $t$-adic order of the error.  The inverse exists
because $(2e-1)^2=1+4(e^2-e)$ and a geometric series in $tB$
converges in $B$.  Completeness gives an idempotent with the
same residue.  But $B\subseteq L$ is a domain and has only
the idempotents zero and one.  Thus $B/tB$, and hence $B$,
has a single maximal ideal $\mathfrak m$.

The ring $B$ is Noetherian, as a finite $A$-module, and
integrally closed in $L$, by transitivity of integrality.
Since $B/tB$ is finite local, $\mathfrak m^r\subseteq tB$
for some least $r\ge1$.  Choose
$y\in\mathfrak m^{r-1}\setminus tB$ and put $v=y/t$.
Then $v\notin B$ and $v\mathfrak m\subseteq B$.
If $v\mathfrak m\subseteq\mathfrak m$, express this action
on finitely many generators of the nonzero ideal $\mathfrak m$.
The determinant trick gives a monic equation for $v$;
the ideal is faithful because $B$ is a domain.  Normality
would give $v\in B$, a contradiction.  Hence the ideal
$v\mathfrak m$ contains a unit and equals $B$.
Consequently $\mathfrak m=\pi B$, where $\pi=v^{-1}\in B$.

The containments $\mathfrak m^{rN}\subseteq t^NB$ make the
$\mathfrak m$-adic topology equivalent to the $t$-adic one.
Every nonzero element therefore has a finite order in powers
of $\pi$, and is a power of $\pi$ times a unit.  Choosing the
least order of a nonzero element in an ideal shows that all
nonzero ideals are powers of $\mathfrak m$.  This proves the asserted
DVR and completeness statements.  Its residue is a quotient
of the finite ring $B/tB$.  Its fraction field is $L$, since
$B$ contains an $E$-basis of $L$.  The inclusion is local
because its inverse image of $\mathfrak m$ is $(t)$.
\hfill$\square$

The standard general counterparts are recorded in the Stacks
Project, tags \href{https://stacks.math.columbia.edu/tag/032L}{032L}
and \href{https://stacks.math.columbia.edu/tag/09E8}{09E8}.
The needed complete equal-characteristic case has been proved
above; those references are not additional assumed producers.
% END CCH9

% BEGIN CCH10
\paragraph{CCH10. The exact power-series presentation and its ramification.}
\label{claim:CCH10}
Let $B$ be as in CCH9, with residue field $k'$ of cardinality
$Q$ and chosen uniformizer $\pi$.  There is a ring isomorphism
\[
                         e:k'[[u]]\xrightarrow{\ \sim\ }B,
                  \qquad e(u)=\pi.                    \tag{CCH.20}
\]
\emph{Proof.}
For each $a\in k'$, start with any lift $b_0$ and iterate
$b_{r+1}=b_r^Q$.  The first difference belongs to $\pi B$;
in characteristic $\ell$, subsequent differences are its
$Q$th powers.  Thus the sequence converges to a root of
$X^Q-X$ lifting $a$.  It is unique: if two such roots differ
by $d\in\pi B$, then $d^Q=d$, and $1-d^{Q-1}$ is a unit,
forcing $d=0$.  Frobenius preserves addition and multiplication,
so these roots form a coefficient field canonically isomorphic
to $k'$.  Denote its section by $[a]$.

Define $e(\sum a_iu^i)=\sum[a_i]\pi^i$.
Completeness gives convergence; multiplication is checked
modulo every $\pi^N$ by finite polynomial multiplication.
A nonzero series has a first nonzero coefficient, whose lift
is a unit, proving injectivity.  Repeatedly subtract the lift
of the residue and divide by $\pi$ to expand any $b\in B$,
proving surjectivity.  These arguments also identify both
adic filtrations.

Let $j=e^{-1}|_A:A\to A'=k'[[u]]$.
It is injective, local and finite, and for some integer
$e_0\ge1$ and some $v\in A'^\times$,
\[
                         j(t)=v u^{e_0}.               \tag{CCH.21}
\]
This follows by the valuation of the nonzero nonunit $j(t)$.
The induced field embedding $\iota_k:k\hookrightarrow k'$
also satisfies
\[
 j(C_k(a))=C_{k'}(\iota_k(a)),\qquad
 \operatorname{const}_{k'}j
          =\iota_k\operatorname{const}_k.              \tag{CCH.22}
\]
For the first equality, the image of $k$ in $B$ is a finite
field mapping injectively into $k'$, so its elements satisfy
$b^Q=b$ and are the unique lifts just constructed.
For the second, write a series as its constant plus $t$ times
another series and use locality.  Thus constants are preserved
through the correct residue-field embedding even when the
uniformizer is ramified.  No assertion $j(t)=u$ is needed.
\hfill$\square$
% END CCH10

% BEGIN CCH11
\paragraph{CCH11. Simultaneous integral and distinct eigenvalues.}
\label{claim:CCH11}
For the finite list $\mathrm{sg}$ from CCH8, let
$P_a(X)=\det(XI-r(\mathrm{sg}(a)))\in A[X]$.
Each is monic of degree $n$ and separable over $E$.
Adjoin all their roots to obtain one finite separable
extension $L/E$.  Separability follows because every adjoined
root satisfies one of the separable $P_a$.
We do not require the product of the $P_a$ to be separable;
different polynomials may share roots.
For $q=0$ take $L=E$.

Every root lies in the integral closure $B$ of CCH9 because
its polynomial is monic over $A$.  Enumerate the $n$ distinct
roots of each polynomial and carry them through (CCH.20)
to elements $\alpha_{a,b}\in A'$, $b\in\mathrm{Fin}(n)$.
Then
\[
 b\longmapsto\alpha_{a,b}\text{ is injective},\qquad
             P_a^j(X)=\prod_{b\in\mathrm{Fin}(n)}(X-\alpha_{a,b}).
 \tag{CCH.23}
\]
The polynomial equality holds in $L[X]$ and hence in $A'[X]$
by injectivity.  These roots need not remain distinct in the
residue field.  The same single finite local map $j$ works
for every selected element.
\hfill$\square$
% END CCH11

% BEGIN CCH12
\paragraph{CCH12. Construct every field of the original coefficient extension.}
\label{claim:CCH12}
The map $j:A\to A'$ of CCH10 gives a point
$x':\texttt{NP}(F,n)$ with the exact original
$\texttt{NP.Ext}(x,x',j)$ and $x'.a=n$.
Give $k'$ its induced $\mathbb Z_\ell$-algebra structure and set
\[
 \rho'=\operatorname{pushRep}(j,\rho),\quad
 \psi'_{v,i}=\operatorname{pushCh}(j,\psi_{v,i}),\quad
 w'_v=\operatorname{GL}(j)(w_v),\quad \mathrm{ty}'=\mathrm{ty}.
 \tag{CCH.24}
\]
The scalar convention is exact:
\[
                 j(\chi(g))=\chi'(g),\qquad j(s)=s',
 \tag{CCH.25}
\]
because the two maps from $\mathbb Z_\ell$ to the
characteristic-$\ell$ ring $A'$ coincide by CCH1.

Here are all the \texttt{NP} fields.  The field $k'$ is finite.
Continuity of each new representation entry follows by
mapping the original quotient modulo $t^N$ to the quotient
modulo $u^N$, using $j((t)^N)\subseteq(u)^N$.
Applying $j$ to (CCH.6) and using (CCH.25) gives \texttt{form}.
The fields \texttt{a\_pos}, \texttt{a\_le} and \texttt{blk}
hold with $a=n$.

For \texttt{span}, choose for each old matrix unit $E_{ij}$
an exponent $b_{ij}$ from $x.\texttt{span}$, and increase
them all to one $b$.  Mapping the finite linear combinations
to $A'$ puts $j(t)^bE_{ij}$ in the span of $r'(H)$.
Multiplication by $v^{-b}$ from (CCH.21) puts
$u^{e_0b}E_{ij}$ there.  Every matrix over $A'$ is an
$A'$-linear combination of these units, proving the full
original \texttt{span} field.

The \texttt{ine} and \texttt{psi\_mu} equalities remain
one after applying the coefficient homomorphism.
The old \texttt{w\_blk} zero entries remain zero with the
same \texttt{ty}.  Mapping the conjugated upper-triangular
matrix gives both parts of \texttt{tri}, with precisely the
same order of the characters.  The two distinctness clauses
in \texttt{N1} are preserved by injectivity of $j$; for the
cyclotomic-twisted comparison use (CCH.25).
This verifies every field listed in
\nolinkurl{PairedDualFirstOrderDualMatrixSpace.lean}:251--274.

By construction, $j$ is injective, local and finite,
$\rho'=\operatorname{pushRep}(j,\rho)$, $x'.a=x.a$, and
all characters are as in (CCH.24).  These are exactly the
six conjuncts of \texttt{NP.Ext}.  Apply CCH2 to $x'$ to
choose its literal $g=0$ Kap, denoted $k'_0$.
\hfill$\square$
% END CCH12

% BEGIN CCH13
\paragraph{CCH13. Transport $N3$, the original $N4$, and the full determinant law.}
\label{claim:CCH13}
For the extension of CCH12, $x.N3$ implies $x'.N3$.
For EVERY subgroup $D\subseteq G$, $x.N4(D)$ implies
$x'.N4(D)$.  Moreover
\[
 \operatorname{GDet.ofRep}(r')
       =\operatorname{GDet.ofRep}(r).\operatorname{map}(j).
 \tag{CCH.26}
\]
\emph{Proof.}
Injectivity of $j$ extends to an embedding of fraction fields.
Characteristic polynomials commute with coefficient maps;
a separable polynomial remains separable after a field
extension, since its Bezout relation with its derivative
remains one.  Thus the same element in each open set
witnessing $x.N3$ also witnesses $x'.N3$.

For $N4$, choose a finite $A$-basis $b_1,\ldots,b_d$ of
$A'$ and its coordinate maps $p_i:A'\to A$.  Such a basis
exists by CCH9 and CCH10.  They satisfy
\[
 y=\sum_i j(p_i(y))b_i,\qquad
                    p_i(j(a)y)=a p_i(y).               \tag{CCH.27}
\]
If $Z'\in M_n(A')$ satisfies
$\chi'(g)r'(g)Z'=Z'r'(g)$ for every $g\in D$,
apply $p_i$ entry by entry.  Equations (CCH.24)--(CCH.25)
and $A$-linearity give
$\chi(g)r(g)p_i(Z')=p_i(Z')r(g)$.
The old $N4(D)$ makes every $p_i(Z')$ zero, and (CCH.27)
makes $Z'$ zero.  This holds for the original auxiliary
$D_0$ passed by CIP10, even though an intermediate frame
record stores $D_0=\bot$.  No replacement of that subgroup
is used in the argument.

For (CCH.26), take any commutative $A'$-algebra $R$ and any
finite group-algebra element $\sum_g a_g[g]$ over $R$.
Both sides evaluate to
\[
                         \det\left(\sum_g a_g r'(g)\right).
\]
Indeed each entry of $r'(g)$ is the coefficient image of
the corresponding entry of $r(g)$.  The determinant is
its integral polynomial in the entries, so this equality
holds over every $R$.  This proves the full polynomial-law
identity, not only its values on individual group elements.

Finally, for any preceding point map $f:B_0\to A$,
injectivity gives $\ker(jf)=\ker f$.
Equation (CCH.22) supplies its exact constant-field and
residual map.  Giving $k'$ the composite action
$O\to k\xrightarrow{\iota_k}k'$ makes $jf$ an
$O$-algebra map whenever $f$ was one.  The determinant and
every ordered-character identity of the original pair
compose with $j$ unchanged.
\hfill$\square$
% END CCH13

% BEGIN CCH14
\paragraph{CCH14. The same finite detector survives the ramified extension.}
\label{claim:CCH14}
For every $S_0\subseteq G$,
\[
 \texttt{Detects}(k_0,S_0)\quad\Longrightarrow\quad
                         \texttt{Detects}(k'_0,S_0).
 \tag{CCH.28}
\]
This is a transport of the same tests; no new elements of
$G'$ are selected after splitting their eigenvalues.

\emph{Proof.}
Use (CCH.27) and $j(t)=vu^{e_0}$.
The coordinate maps have the exact continuity property
\[
               p_i(u^{e_0N}A')\subseteq t^NA.
 \tag{CCH.29}
\]
Indeed $u^{e_0N}A'=j(t^N)A'$, and apply $A$-linearity.
Take an $A'$-valued polarized twisted cocycle $z'$ whose
trace tests vanish on $S_0$, as in CCH3 for $x'$.
Let $z_i=p_i(z')$ entrywise.  Equation (CCH.29) proves
the required adic continuity of $z_i$.  Applying $p_i$ to
the cocycle and transpose equations gives (CCH.8) over $A$,
because all the matrices $r'(g)$ and scalars $\chi'(g),s'$
come from $A$.  Likewise
\[
 \operatorname{tr}(z_i(h)r(h))
                  =p_i\bigl(\operatorname{tr}(z'(h)r'(h))\bigr)=0
\]
on $S_0$.  The old detector gives integers $a_i\ge0$ and
matrices $Y_i\in M_n(A)$ with
$sY_i^T=Y_i$ and $t^{a_i}z_i=d_\chi Y_i$.

Choose $a\ge a_i$ for all the finitely many coordinates and set
\[
 Y'=v^{-a}\sum_i b_i\,j(t^{a-a_i}Y_i)\in M_n(A').
 \tag{CCH.30}
\]
The maps on matrices are entrywise.  Equations
(CCH.25) and (CCH.27) give $s'(Y')^T=Y'$ and
\[
 \begin{split}
 d_{\chi'}Y'
 &=v^{-a}\sum_i b_i j(t^{a-a_i}d_\chi Y_i)\\
 &=v^{-a}j(t)^a\sum_i b_i j(z_i)
   =u^{e_0a}z'.
 \end{split}                                           \tag{CCH.31}
\]
CCH3 identifies this with the literal \texttt{TwB} conclusion.
The exponent $e_0a$ is allowed to depend on the tested cocycle;
the list $S_0$ and the coefficient extension were already fixed.
\hfill$\square$
% END CCH14

% BEGIN CCH15
\paragraph{CCH15. The complete original one-point \texttt{Choice} output.}
\label{claim:CCH15}
Under (CCH.2), there exist
\[
 \begin{gathered}
 q,\quad \mathrm{sg}:\mathrm{Fin}(q)\to G,\quad
 x':\texttt{NP}(F,n),\quad j:A\to x'.A,\\
 k'_0:\texttt{Kap}(x'),\quad
 \alpha:\mathrm{Fin}(q)\to\mathrm{Fin}(n)\to x'.A
 \end{gathered}
\]
such that
\[
 \begin{gathered}
 \mathrm{sg}(a)\in F.G',\qquad x.\texttt{Ext}(x',j),\\
 k'_0.g=0,\qquad \texttt{Detects}(k'_0,\operatorname{range}\mathrm{sg}),\\
 \alpha_a\text{ is injective},\qquad
 (\operatorname{mat}x'.\rho(\mathrm{sg}(a))).\mathrm{charpoly}
                  =\prod_b(X-\alpha_{a,b}).
 \end{gathered}
 \tag{CCH.32}
\]
To construct them, choose $k_0$ by CCH2 and the finite list by
CCH8.  Apply CCH9--CCH11 to its characteristic polynomials.
Construct $x'$ and $k'_0$ by CCH12, and use CCH14 on the
same list.  Formula (CCH.23), together with coefficientwise
matrix and characteristic-polynomial transport, gives the
last line of (CCH.32).  CCH13 additionally gives $N3$, every
required $N4(D)$, the full determinant identity and kernel
preservation on this same extension.

For the source's $h=0$ family indexed by $\mathrm{Fin}(1)$,
assign $x',j,k'_0,\alpha$ to its sole index.  The remaining
residue-cardinality requirement compares this field with
itself, and is reflexive.  Thus every conjunct of
\nolinkurl{PatchedSupportCompletions.lean}:223--230 is
supplied at the original one-point full-block scope.
This is exactly the scope supplied by CIP10 before the
\texttt{NiceI} consumer; a multi-point or arbitrary-block
\texttt{Choice} theorem is not required here.
\hfill$\square$
% END CCH15

% BEGIN CCH16
\paragraph{CCH16. Literal $B2$ for this chosen Kap, with the full level quantifiers.}
\label{claim:CCH16}
Let $T$ be any original \texttt{TWGrp} on the relevant
\texttt{Glob}, and let $x'$ be the point just constructed.
For its chosen $k'_0$ the literal $T.B2(k'_0)$ holds with
the single constant $C=0$.

\emph{Proof.}
Choose $C=0$ before the truncation $m$.  For every $m$ the
conclusion in fact holds at every level $N$, which implies
the required eventual statement.
The source quantifies over an array
\[
 z:\mathrm{Fin}(k'_0.g+1)\to\mathrm{Fin}(k'_0.g)
       \to(T.G_N)_0\to M_n(x'.A/(u^m)).
\]
Since $k'_0.g=0$, its second index set is empty.  Choose
$c:\mathrm{Fin}(1)\to x'.A$ to be constantly one; its value
at $0$ is a unit.  Choose the zero function $\Psi$ in
\texttt{Sel2}.  To check this membership directly, zero
lies in the coefficient submodule $W2$; all quotient maps
of zero are locally constant; the cocycle, conjugation and
tame-vanishing equations are zero equations; and each local
\texttt{LocAt} condition holds with the zero correction
matrix.  This also handles $m=0$, when the quotient ring is zero.
The final required equality of the $\mathrm{ep}'_j$ coordinates
is indexed by $j\in\mathrm{Fin}(0)$ and is vacuous.

Thus $\exists C\,\forall m\,\forall^{\mathrm{eventual}}N$
has been proved with $C=0$ and with every $N$ allowed.
The proof is for the actual chosen $g=0$ Kap.  It does not
assert the source's universal obstruction theorem for an
arbitrary Kap with a nonempty odd Selmer family.
\hfill$\square$
% END CCH16

% BEGIN CCH17
\paragraph{CCH17. Install the coupled outputs in the original same-ideal consumer.}
\label{claim:CCH17}
Fix the original rank-four data, the same universal pair
$(B,\Lambda)$ and an ideal $J\subseteq B$ already selected
by the ordinary construction.  Retain original \texttt{CohP},
\texttt{LocCFT}, \texttt{IsCoef}, $V_0.S'=s.S$, the
nontrivial residual cyclotomic value, the same-ideal
\texttt{HasAut}, the original first-order bound, and the
one-point \texttt{Raw} support conclusion.
Then the \texttt{Choice} and obstruction-bound inputs to
the \texttt{NiceI} part of JSP11 are both supplied by
CCH15--CCH16 on one common chosen point and Kap.

\emph{Proof.}
Take an arbitrary original \texttt{P1core} point and its map
$f:B\to x_0.A$.  CIP10--CIP11 normalize it on the original
updated inertia frame, producing $x.a=4$, $x.N3$ and
$x.N4(D_0)$ for the old auxiliary subgroup, with an
injective local map and all determinant/character identities.
Take the $\Phi_\ell$ part of CAP's \texttt{ArPl}, under the
unchanged \texttt{CohP}.  Apply CCH15 to obtain the finite
list, the second extension and its exact $g=0$ Kap.
CCH13 composes the point map, preserves its kernel and the
original $N4(D_0)$, and supplies the required $O$-algebra
and residual constant-field convention.

CAP13 realizes precisely this list by the original $T$
and its actual \texttt{TWPl}.  Its equality for the residual
kernel $H$ and the full residual determinant law give the
same PL-kernel condition already proved in JSP11.  The
retained first-order theorem therefore supplies
\texttt{B1} with the exact parameter
\[
                       4[L^+:\mathbb Q]+4T.q.
\]
CCH16 supplies \texttt{B2} for this SAME $k'_0$ and this
SAME $T$, with no new arithmetic bound assumed.
The eigenvalues, detector, point map and both bound fields
now form the literal original \texttt{Pt}.
JSP9--JSP11 supply the unchanged \texttt{Adapted},
\texttt{GenBy}, stable-discriminant and kernel identifications.
The retained \texttt{HasAut}(B,J) provides its system over
this $T/\texttt{TWPl}$ on the already fixed ideal $J$.
Applying the retained $h=0$ \texttt{Raw} conclusion and
using the unchanged point kernel gives the original
\texttt{P1core} statement, as in JSP11.

All choices occur inside that point quantifier.  They do
not reselect $J$, enlarge the external coefficient object,
or change the earlier threshold.  In the two-Link Count-II
application, fix $\iota$ first, then choose the common
coefficient threshold for the two original rank-four Links,
then allow every later coefficient and arbitrary actual
$B,\Lambda$.  Each Link has its own ring and one ideal
carrying all six Hida fields and \texttt{P1core}; no equality
between ideals in different rings is asserted.
The present proof works at $\ell=3$ and does not invert $24$.
\hfill$\square$

The selected one-point \texttt{Choice} producer is now
constructed.  The selected obstruction input is discharged
through its coupled $g=0$ choice, while the general source
obstruction theorem remains outside the claim.
The five retained Count-II producer entries are original
ordinary/control/arithmetic-fiber data (including integral
type, \texttt{neatU} and \texttt{TThetaEx}), four further
Hida endpoints, same-$J$ all-level \texttt{HasAut}, the
first-order bound, and one-point \texttt{Raw}.
\texttt{CohP} and all stated side conditions remain.
QX10 and YX2 keep the same route from the joint two-Link
supplier to \texttt{CountII}; the whole supplier is still open.
KA5--KA6 remove only Case III for two-primary targets.
Cases I, II and V still contribute to RH, and unrestricted
Dirichlet GRH retains its additional Case-III/SeedIV chain.
% END CCH17
